\uuid{Q5t8}
\titre{ Racine carrée d'une matrice }

\niveau{L2}
\module{Algèbre}
\chapitre{Réduction d'endomorphisme}
\sousChapitre{Applications}
\cpge{mpsi-mp2i.calcul-matriciel-systemes.e, pcsi.determinants.c, ptsi.determinants.a, mpsi-mp2i.changements-bases.c, bcpst1.matrices, mp-mpi.reduction-endomorphismes.c, pc.reduction.b, psi.reduction-endomorphismes.b, pt.reduction-endomorphismes.b, bcpst2.valeurs-propres-vecteurs-propres.a, bcpst2.valeurs-propres-vecteurs-propres.b}
\theme{}
\auteur{Quentin Liard}
\datecreate{ 2026-03-05}
\organisation{AMSCC}
\difficulte{1}

\contenu{

\texte{
On considère la matrice
\[
A=
\begin{pmatrix}
4 & 0\\
0 & 9
\end{pmatrix}.
\]
}

\begin{enumerate}

\item \question{Déterminer les valeurs propres de $A$.}
\reponse{$4$ et $9$.}

\item \question{Montrer que $A$ est diagonalisable.}
\reponse{Matrice diagonale donc diagonalisable.}

\item \question{Déterminer une matrice $B$ telle que $B^2=A$.}
\reponse{
\[
B=
\begin{pmatrix}
2 & 0\\
0 & 3
\end{pmatrix}.
\]
}

\end{enumerate}

}